\section{Separate class dimensions in the triangular cut}
\label{sec:cubic-sharpened-dimension-gap}

This section isolates the numerical estimate actually used in the
triangular part of the source dimension argument.  The two adjacent
cubic roots are conjugate, so their coordinates require one strict
class space, not two independent copies.  The central class space
and the trivial-character class space have separate roles.  No
vanishing of a cubic Iwasawa mu invariant is assumed.

\paragraph{Actual tuple and algebra.}
Let $p=269$, let $F$ be one of the totally real selected-place hosts
of AC-G93-fill, and put $H=FK$, where $K=\mathbb Q(\omega)$.
Write $\Gamma=G_F/N$ for the exact residual-away-kernel quotient
of AC-G90-quotient, and $G=\ker(\eta_K:\Gamma\to\{\pm1\})$.
Complex conjugation $c\in\Gamma$ has order two.  The local maps are
all full $G_{F_v}$, $v\mid p$, and their two lifts to $G$; the
extension $H/F$ splits at these places.  Put
\[
 d=\sum_{v\mid p}d_v=[F:\mathbb Q],\qquad
 d_v=[F_v:\mathbb Q_p],\qquad
 \theta_0=\chi^j\gamma,\quad\theta_1=1,\quad
 \theta_2=\chi^j\gamma^{-1},\quad j\in\{1,2\}.
 \tag{G98.1}
\]
Here $\chi^c=\chi^{-1}$ and $\gamma^c=\gamma$.
All characters take values in a finite field $k$ of characteristic
$p$.  Set
\[
 \alpha=\theta_0^{-1},\qquad \beta=\theta_2=\alpha^c,
 \qquad\omega=\alpha\beta=\gamma^{-2},\qquad
 W=\gamma^{-2}\eta_K.
 \tag{G98.2}
\]
The character $W$ is defined on $\Gamma$; the two adjacent
characters need only be defined on $G$.
Retain the actual local separation and normalization of AC-G94-space:
at each local map the three residual characters have distinct values
on some element, and at a fixed $g_*\in G$,
$\theta_2(g_*)=\theta_0(g_*)^{-1}$ and
$\theta_0(g_*)^2\ne1$.  The universal representation satisfies
\[
 r(g_*)-r(cg_*^{-1}c)\ \hbox{is diagonal}.
 \tag{G98.3}
\]
For the actual tuple one may use the cubed local separating element
of AC-G90-Low; the cubic character then has value one.

Let $R,B,A,x_0=\mathfrak m_BA$ be the actual rings of AC-G95.
In particular $B=R_0$ is complete Noetherian local, $A$ is finite
type and graded over $B$, and its generators are actual matrix
entries, with weights $w_i-w_j$ for $w=(1,0,-1)$.
All diagonal entries and local characters lie in $B$.
Modulo $x_0$ they equal their prescribed residual values.
Entries are locally constant modulo $x_0$.
Define
\[
 T_-=(r_{ij}(g):i<j),\qquad T_+=(r_{ij}(g):i>j)
           \quad\hbox{in }A.
 \tag{G98.4}
\]
Thus $T_-=0$ means lower triangular; it is the source's
\texttt{triM}.  The ideal $T_+$ is its \texttt{triP}.
Put
\[
 \begin{split}
 a&=\dim_k H^1_{\mathrm{str}}(G,k(\alpha)),\\
 z&=\dim_k H^1_{\mathrm{str}}(\Gamma,k(W)),\\
 h&=\dim_k H^1(\Gamma,k).
 \end{split}
 \tag{G98.5}
\]
The strict conditions use every full local map.  Local separation
makes the degree-one cone equal to the strict kernel for $\alpha$
and $W$.  These dimensions are finite by AC-G74, AC-G76 and the
exact quotient comparison in AC-G90; no numerical upper bounds
are imposed in their definition.

\paragraph{Local arithmetic hypotheses used by the cut.}
We use that every $F_v$ contains no primitive $p$th root of unity,
and that additive $p$-power characters of $G$ unramified at all
finite places have bounded exponent.  These facts hold on the
actual hosts here.  At the initial quadratic stage the $p$-completions
are unramified extensions of $\mathbb Q_p$ of degree one or two;
the G83 cover adds only unramified $p$-power local degree, the G92
selected-place extension splits at $p$, and a subsequent
cyclotomic layer has $p$-power local degree.  Thus a local degree
is prime to $67$, whereas $[\mathbb Q_p(\mu_p):\mathbb Q_p]=268$.
The bounded exponent is supplied by the finite ideal class group
of $H$.  An additive $p$-group character factoring through the
residual-away quotient is already unramified away from $p$:
its allowed inertia quotient has order prime to $p$.
These are qualitative facts, not a class-number bound.

\begin{theorem}[AC-G98-dimension: the actual separate-dimension cut]
\label{thm:ACG98dimension}
The lower residual fiber satisfies
\[
 \dim A/(x_0+T_-)\le a+z.                         \tag{G98.6}
\]
Let $x$ be a closed point of $\operatorname{Spec}A$ above
$\mathfrak m_B$, with $T_+\not\subset x$.
Let $C$ be the completed local ring at a point above $x$ after
an arbitrary finite unramified coefficient extension of
$\mathcal O$.  This is the source's actual local-completion
construction; its residue field may be a finite extension of $k$.
Then
\[
 \dim C/\mathfrak m_BC\le a+z,\qquad
 \dim C/T_-C\le 2d+h+a+z+1.                     \tag{G98.7}
\]
Consequently the triangular exclusion needed for the rank-three
source argument follows from
\[
                         h+a+z+1<d.             \tag{G98.8}
\]
In the old one-character notation the literal source hypotheses
only give $h\le1+C_0$, $a\le C_0$, $z\le C_0$.
The proof therefore gives $2d+3C_0+2$ there.  Its stated
$2d+3C_0+7$ is a looser bound, not an additional contribution
from a cohomology group.

The theorem proves this triangular cut for the actual cubic
tuple.  It does not assert every other source dimension or
propagation theorem for the changed tuple.  In particular the
local character-collision argument, the local presentation
bounds, arithmetic counting and a prescribed Hecke ideal remain
separate interfaces.  If the retained collision argument is used,
its existing condition $d_v>5$ may be kept unchanged.
\end{theorem}

\begin{proof}
\emph{1. Normalized adjacent coordinates.}
Let $Q$ be any commutative $k$-algebra and let $r$ be a lower
triangular point of the actual residual fiber with locally constant
entries.  The diagonal characters are exactly
$(\theta_0,1,\theta_2)$.  Write
\[
 r(g)=U(x(g),y(g),t(g))D(g),\qquad
 D=\operatorname{diag}(\theta_0,1,\theta_2).
 \tag{G98.9}
\]
The complete matrix computations of AC-G90 give
\[
 d_\alpha x=0,\quad y=-x^c,\quad
 d_\omega t=-y\smile x,\quad t+t^c=xy.
 \tag{G98.10}
\]
The local ordered identity is exactly constant local splitting,
including over $Q$, since the separating differences are units.
Thus $x$ is strict, with unique local primitives.  Its normalization
is the linear equation obtained from the $(1,0)$ entry of (G98.3).
A coboundary $x(g)=(\alpha(g)-1)u$ gives
$r_{10}(g)=(1-\theta_0(g))u$.  Its normalization difference is
\[
                    (\theta_2(g_*)-\theta_0(g_*))u.
 \tag{G98.11}
\]
The coefficient is nonzero in $k$.  Hence each strict cohomology
class has a unique normalized representative, by subtracting
one coboundary.  No residual absolute irreducibility is needed.

Continuous cochains with values in a discrete $k$-vector space
$V$ equal the tensor product of the $k$-cochains with $V$:
every continuous map from a profinite space has finite image,
which spans a finite-dimensional subspace of $V$.
The same holds for the finite sum of local cochains.
Tensoring over the field $k$ is exact, so it commutes with the
strict cohomology, the normalization and the involution projectors.
Choose a basis of the normalized strict adjacent space over $k$.
There are at most $a$ coefficients $q_1,\ldots,q_a\in Q$
such that every value of $x$, hence of $y$, belongs to
\[
                         U=k[q_1,\ldots,q_a]\subset Q.
 \tag{G98.12}
\]
Each local primitive is in $U$, since it is an evaluation of $x$
divided by a fixed nonzero local character difference.

\emph{2. The central coordinate after the adjacent algebra.}
Pass to the $k$-vector-space quotient $V=Q/U$, not a ring quotient.
The cup term and $xy$ in (G98.10) disappear.  The image $\bar t$
is a strict $\omega$-cocycle satisfying
$\bar t^c=-\bar t$; the local primitives have the same anti-symmetry,
by AC-G90.17.  Quadratic descent by the exact projector
$(1-c)/2$ identifies its cohomology with
$H^1_{\mathrm{str}}(\Gamma,k(W))\otimes_kV$.
There is no additional central coboundary freedom in this
anti-space.  A boundary $d_\omega u$ is invariant under
precomposition by $c$, since $\omega^c=\omega$.
Being also anti-invariant forces it to vanish.  Equivalently,
a cocycle on $\Gamma$ with coefficient $W$ has a unique representative
with value zero at $c$, because $W(c)=-1$; its restriction is anti.
Local invariants vanish, so there is no extra local-primitive kernel.

Choose a basis of this central space, and lift its at most $z$
coefficients in $Q/U$ to $q'_1,\ldots,q'_z\in Q$.
Every value of $t$ belongs to $U[q'_1,\ldots,q'_z]$.
Apply this to $Q=A/(x_0+T_-)$.
This $k$-algebra is generated by the displayed matrix entries,
by AC-G95; its local characters and diagonal entries are constants.
It is therefore generated by at most $a+z$ elements.
It is a quotient of a polynomial algebra in that many variables,
which proves (G98.6).  This argument neither assumes that every
adjacent class has a filling nor counts the conjugate root twice.

\emph{3. The complete special fiber at a nonupper point.}
In any field-valued point of $A/x_0$, a nonzero entry of positive
weight and a nonzero entry of negative weight cannot coexist.
Indeed take positive integers making their two weights sum to zero.
The resulting monomial is in $B\cap\mathfrak m_R=\mathfrak m_B$,
because it contains an off-diagonal factor.  Its image must be zero,
contradicting the field property.  Thus every prime of $A$ containing
$x_0$ contains $T_+$ or $T_-$.
At primes lying below the specified $x$, the first alternative is
impossible.  All these primes contain $T_-$.
Consequently the image of $T_-$ in $A_x/x_0A_x$ is nilpotent,
and this local special fiber has dimension at most $a+z$ by (G98.6).
Here is the exact completion argument, with no appeal to the
source's one-character declaration.  Write
$D=(A\otimes_{\mathcal O}\mathcal O')_y$, where $y$ lies over
$x$ and $\mathcal O'/\mathcal O$ is the stated finite unramified
extension.  The map $A_x\to D$ is an etale local homomorphism
of Noetherian rings, as is its base change modulo $x_0$.
Its closed fiber has dimension zero.  The flat local dimension
formula gives
$\dim(D/x_0D)=\dim(A_x/x_0A_x)$.
Since $C=\widehat D$, exactness of completion on finite modules
gives $C/\mathfrak m_BC\simeq\widehat{D/x_0D}$.
The dimension of a Noetherian local ring equals that of its
completion.  These are respectively Stacks Lemmas~41.19.1 and
10.112.7, Lemma~10.97.1, and Lemma~15.44.1; all their Noetherian,
local, finite-module and etale hypotheses have just been verified.
A finite residue-field extension also preserves $a,z,h$ when
dimensions are taken over the enlarged residue field, by exact
scalar extension of the cochain complexes.
This proves the first bound in (G98.7).

\emph{4. Exactly $h+d$ tests for an endpoint character.}
First choose $h$ elements $t_i\in\Gamma$ on which evaluation
is injective for $H^1(\Gamma,k)=\operatorname{Hom}_{\rm cont}(\Gamma,k)$.
This is possible by finite-dimensional linear algebra.
They detect every character into a finite abelian $p$-group:
a nontrivial image admits a nonzero quotient to the additive
prime field, which would contradict that injectivity.

Local class field theory identifies the maximal pro-$p$ quotient
of $\mathcal O_{F_v}^\times$ with $\mathbb Z_p^{d_v}$, since
$\mu_p\not\subset F_v$.  Choose $d_v$ inertia elements generating
it.  Raise them to the prime-to-$p$ exponent $\#k^\times$;
they still generate the same pro-$p$ quotient, and all residual
characters are now trivial on them.
Their images in $G$ give $d$ elements $s'_j$.
For any finite abelian $p$-group character of $G$ invariant under
$g\mapsto cg^{-1}c$, vanishing on these $s'_j$ means vanishing
on the inertia at both members of every local pair.
It is unramified everywhere.  A fixed nonzero integer $m$,
for example the exponent of the $p$-class group of $H$,
therefore kills its image.

Write any element of $\Gamma$ as $u$ or $uc$, with $u\in G$;
denote this set-theoretic retraction by $\pi$.
Adjoin to the $s'_j$ the $h$ elements
\[
 s_i=(\pi(t_i)c\pi(t_i)c)^{\#k^\times}.
 \tag{G98.13}
\]
If a finite abelian $p$-group character $\nu$ of $G$ kills
all these $h+d$ elements, then $\nu\nu^c$ descends to
$\Gamma$ with value one at $c$ and vanishes on every $t_i$.
The prime-to-$p$ exponent is invertible on its finite image.
The preceding detection proves $\nu^c=\nu^{-1}$.
It now kills both sets of local inertia, and $\nu^m=1$.
Thus these $h+d$ elements control a general endpoint character;
only the $d$ elements $s'_j$ are needed for the self-dual middle
character.  This reproduces the complete \texttt{AbEval} proof
with $h$ in place of the uniform bound $1+C_0$.

\emph{5. The triangular cut and its count.}
In $C/T_-C$, use the following elements:
\[
 \varpi,\quad r_{00}(s_i)-1\ (h+d\hbox{ elements}),\quad
 r_{11}(s'_j)-1\ (d\hbox{ elements}).
 \tag{G98.14}
\]
There are exactly $2d+h+1$ of them, all in the maximal ideal.
Consider a prime $P$ containing $T_-C$ and all these elements.
The complete local domain $C/P$ has characteristic $p$.
Its diagonal entries are multiplicative characters; divide them
by their prescribed residual characters.  The resulting characters
take values in $1+\mathfrak m_{C/P}$ and are adically continuous.
Each reduction modulo a power of the maximal ideal has finite
$p$-group image.  Step 4 implies that the normalized first and
middle characters have exponent dividing $m$ in every reduction,
hence exactly in $C/P$.  They are trivial: in a characteristic-$p$
local domain, the only finite-order element of $1+\mathfrak m$
is one.  To see this, remove the $p$-power part using
$X^{p^r}-1=(X-1)^{p^r}$, and for the prime-to-$p$ part the
cofactor of $X-1$ is a unit.  Polarization gives the third diagonal
character as well.

The local ordered character values are also their prescribed
constants.  The characteristic polynomial identity says that
each is a root of the product of the three diagonal linear factors.
In the domain it equals one of their $k$-valued roots.  Reduction
then selects its prescribed value, because $k$ embeds in the residue
field.  A zero-weight monomial containing an off-diagonal entry
is zero in a lower triangular matrix family.  It follows that
all defining generators of $B$ have their residual values.
The same holds for their adic closure: a finite residue field is
closed in $C/P$, and the maps are continuous.
Thus $\mathfrak m_BC\subset P$.

The radical of the ideal obtained by adding (G98.14) to $T_-C$
therefore contains $\mathfrak m_BC$.
The generalized principal ideal dimension inequality in the
Noetherian local ring $C/T_-C$ gives
\[
 \dim C/T_-C
 \le (2d+h+1)+\dim C/(T_-C,(\mathrm{G98.14}))
 \le 2d+h+1+a+z.
 \tag{G98.15}
\]
This is exactly the source's \texttt{KrullDimQuotientBound}
argument, now with each number of generators displayed.
It proves the second bound in (G98.7) and (G98.8).
The substitution $h\le1+C_0$, $a,z\le C_0$ proves the
comparison with the original uniform estimate.
\end{proof}

\begin{theorem}[AC-G98-slope: a finite adjacent certificate suffices]
\label{thm:ACG98slope}
Retain the actual selected-place host of AC-G93-fill and its
cyclotomic layers $F_n$, $H_n=HF_n$.  Begin after eventual total
ramification at all $p$-places.  Then $d_n=d_0p^n$.
Let $a_n,z_n,h_n$ be (G98.5) on this tower.
There are a nonnegative integer $r_\alpha$ and finitely many
positive integers $t_i$ such that
\[
 a_n=r_\alpha p^n+\sum_i\min(t_i,p^n),\qquad
                   h_n=O(1),\quad z_n=O(1).
 \tag{G98.16}
\]
Consequently the sharper numerical condition
\[
                     h_n+a_n+z_n+1<d_n          \tag{G98.17}
\]
holds at every sufficiently large layer if and only if
\[
                             r_\alpha<d_0.      \tag{G98.18}
\]
In particular one finite strict adjacent calculation
\[
                            a_m<d_0p^m          \tag{G98.19}
\]
is sufficient; conversely such a finite level exists whenever
(G98.18) holds.  This allows a positive cubic growth rank.
Neither (G98.18) nor (G98.19) is established for the actual cubic
field by this theorem.  They replace the stronger common-bound
slope test $3\max(R_1,R_2)<d_0$ for this triangular dimension cut.
They do not by themselves establish the full literal
\texttt{Frame.Big} declaration or transport the remaining source
dimension and Hecke-containment statements.
\end{theorem}

\begin{proof}
\emph{1. The exact adjacent sector and no permutation correction.}
Take the minimal finite extension $L/H$ splitting $\alpha$,
with group $\Delta$, and put $L_n=LH_n$.
Its degree is prime to $p$.  The exact away-kernel comparison
of AC-G90 and AC-G74 identifies
\[
 a_n=\dim_k e_\alpha
       \bigl(\operatorname{Cl}(\mathcal O_{L_n}[1/p])/p
                                  \otimes_{\mathbb F_p}k\bigr).
 \tag{G98.20}
\]
The character orientation is $\alpha$, not $\alpha^{-1}$.
Let $X'$ be the maximal unramified, $p$-split Iwasawa module
over $L_\infty$, and let
$M_\alpha=e_\alpha((X'/p)\otimes_{\mathbb F_p}k)$.
AC-G74 proves that it is a finitely generated $k[[T]]$-module and gives
\[
 \mathbb Z_p[\Sigma_p]^0\longrightarrow
       (X')_{\Gamma_{\mathrm{cyc}}^{p^n}}
       \longrightarrow
       \operatorname{Cl}(\mathcal O_{L_n}[1/p])_{(p)}
       \longrightarrow0.
 \tag{G98.21}
\]
Every decomposition subgroup in $\Delta$ at $p$ has a
nontrivial $\alpha$-restriction.  Indeed the local cubic
inertia has order three, whereas the image of $\gamma$ has
order dividing $\#k^\times$ and, on the actual coefficient
inertia, divides $268$.  These local orders are coprime.
The initial host is unramified at $p$; selected-place extensions
split there and the further extensions have $p$-power degree.
They cannot kill the order-three character.
Thus the cubic factor in $\alpha$ cannot cancel at any $p$-place.

Frobenius reciprocity for the permutation module gives
$e_\alpha k[\Sigma_p]=0$, hence also
$e_\alpha k[\Sigma_p]^0=0$.
Reducing (G98.21) modulo $p$ and applying the exact idempotent
therefore gives the actual isomorphism
\[
 e_\alpha\bigl(\operatorname{Cl}(\mathcal O_{L_n}[1/p])/p
                                     \otimes k\bigr)
                 \simeq M_\alpha/T^{p^n}M_\alpha.
 \tag{G98.22}
\]
The structure theorem over the discrete valuation ring $k[[T]]$
now gives the first formula in (G98.16), with
$r_\alpha=\operatorname{rank}_{k[[T]]}M_\alpha$.
No vanishing or numerical bound on this rank has been used.

\emph{2. Only rational coefficient fields enter the bounded terms.}
The field $F$ is obtained from a real abelian field by the fixed
finite $p$-power host constructions and the G92 selected-place
extension.  Its normal closure over a suitable real abelian
field $A_0$ is a finite $p$-extension, as checked in AC-G93-fill.
Let $B_W/\mathbb Q$ be the abelian field cut out by the rational
character $W=\gamma^{-2}\eta_K$.
The fields $A_0(\mu_p)$ and $A_0B_W(\mu_p)$ are abelian
over $\mathbb Q$.
Ferrero--Washington gives cyclotomic $\mu=0$ for them.
The finite-$p$ ascent proved in AC-G76-host, applied through a
normal $p$-closure and a subgroup chain, gives the same assertion
for $F(\mu_p)$ and $FB_W(\mu_p)$.
This also applies when the selected-place host itself is not
Galois over $\mathbb Q$: each successive index-$p$ subgroup
in the finite $p$-group is normal in its predecessor, giving
cyclic degree-$p$ field steps inside the desired field.

Their unramified and $p$-split Iwasawa modules are finitely
generated $\mathbb Z_p$-modules.  By AC-G74, their finite-layer $p$-split
class ranks are bounded.  Exact prime-to-$p$ descent identifies
the central strict group with the $W$-sector in the corresponding
class group, so $z_n=O(1)$.
For the trivial full-local group, AC-G76.2 gives
\[
 h_n=1+a_{F_n}(\bar\epsilon)+u_{F_n}.
 \tag{G98.23}
\]
The class term is bounded by the first rational coefficient tower,
and $u_{F_n}$ is bounded by the fixed number of $p$-places
(in fact it is zero for the stated local fields).
This proves the remaining bounds in (G98.16).
At no point is Ferrero--Washington applied to a field needed
to split the cubic character $\chi$.

\emph{3. The precise slope and the finite test.}
The torsion sum in (G98.16) is nonnegative and bounded.
If $r_\alpha<d_0$, the difference
$(d_0-r_\alpha)p^n$ eventually exceeds that sum plus the
bounded $h_n+z_n+1$, proving (G98.17).
If $r_\alpha\ge d_0$, then $a_n\ge d_0p^n=d_n$,
so (G98.17) is impossible at every layer.
Finally $r_\alpha p^m\le a_m$ proves the sufficiency of
(G98.19); the converse follows by taking $m$ large.
All local degrees $d_{v,n}$ also eventually exceed five,
so the retained local degree requirement costs no further
asymptotic class bound.  This last observation is not a proof
of a changed-tuple local presentation or propagation theorem.
\end{proof}

\paragraph{Exact source proof interface.}
All Lean sources in this paragraph are pinned to commit
\texttt{7cbf7a0c32fca1badf05931879a37c3f8a2b03b4} of
\url{https://github.com/CocoSgt/riemann-hypothesis-lean}, below
\path{src/Development/AutomorphicArithmetic/AutomorphicVarieties/FrobeniusSmoothVectors/}.
The complete applied proofs freshly read are
\path{AbEval.lean}, \path{FinrankStructMapYRange.lean},
\path{FinrankStructMapSRange.lean},
\path{FramePointYCoordEq.lean}, \path{FramePointAlgCoordEq.lean},
\path{FinrankFrameSpaceEq.lean}, \path{FibDim.lean},
\path{IrrOfConn.lean}, \path{CharpolyCoeffsInSubalgebra.lean},
\path{TriangularRepresentationDiagonal.lean},
\path{TrdegQuotLeOne.lean}, \path{LocDim.lean},
\path{KrullDimQuotientBound.lean},
\path{TameInertiaCondition.lean}, and
\path{KrullDimLtOfNotLeTriangular.lean}.
The source's uniform $2C_0$ lower-fiber bound adds the adjacent
and central coordinate counts; the present proof supplies both
maps for the actual residual tuple instead of substituting it
into the source's literal one-character declarations.
The complete applied Stacks proofs are Lemma~15.44.1,
\url{https://stacks.math.columbia.edu/tag/07NV};
Lemma~41.19.1 in
\url{https://stacks.math.columbia.edu/tag/025L}, with its cited
dimension proof Lemma~10.112.7,
\url{https://stacks.math.columbia.edu/tag/00ON}; and
Lemma~10.97.1 in
\url{https://stacks.math.columbia.edu/tag/0BNH}.
Their earlier Hilbert--Samuel, flat going-down and Artin--Rees
inputs remain imported, as does the generalized principal ideal
dimension inequality used in the fully read source proof.
No Lean execution,
new source automorphy proof, or unproved cubic mu vanishing
is claimed.

\paragraph{Geometric dictionary and remaining arithmetic gate.}
The closed lower fiber is a finite-type moduli space of normalized
extensions of the three fixed local systems.  The first adjacent
extension determines the conjugate adjacent extension by polarity.
Modulo the algebra it generates, the remaining coordinate is a
central strict extension.  The cut counts generators of these
spaces and diagonal characters separately.
The class-sector description uses the already checked finite-etale
cover and character-projector dictionary of SGA~1, Expose~V;
no new archival statement is attributed to Grothendieck here.
The outstanding arithmetic task is an actual certificate
$a_m<d_m$ on a suitable retained host, or another proof of
$r_\alpha<d_0$.  Genus positivity and a proof of the Euler identity
supply no such upper estimate.
